
\begin{derivation}[从\eqnrefs{eq:phi4-bubble-integral-r68}到\eqnrefs{eq:phi4-bubble-result-dp10}]
\eqnrefs{eq:dim-reg-measure-general-dp10} 中显式 \(\hbar\) 的消去可由变量替换直接看出。若从逆长度圈变量 \(\kappa^\mu\) 开始，定义物理动量

\begin{equation*}
 \ell^\mu=\hbar\kappa^\mu,
 \qquad
 \mu_{\rm DR}=\hbar\mu_\kappa,
\end{equation*}
则
\begin{equation*}
 \dd^d \kappa=\frac{\dd^d \ell}{\hbar^d},
 \qquad
 (\kappa^2-M^2)^{-2}
 =\frac{\hbar^4}{(\ell^2-m_R^2c^2)^2},
\end{equation*}
其中 \(M=m_Rc/\hbar\)。维数正则化尺度给
\begin{equation*}
 \mu_\kappa^{2\epsilon}
 =\frac{\mu_{\rm DR}^{2\epsilon}}{\hbar^{2\epsilon}}.
\end{equation*}
由于 \(d=4-2\epsilon\)，总的 \(\hbar\) 次数为
\begin{equation*}
 4-d-2\epsilon=0.
\end{equation*}
因此变量换元后的物理四动量形式为
\begin{equation*}
 \mu_{\rm DR}^{2\epsilon}
 \int\frac{\dd^d \ell}{(2\pi)^d}
 \frac{\ii}{(\ell^2-m_R^2c^2+\ii0^+)
 ((\ell+p)^2-m_R^2c^2+\ii0^+)},
\end{equation*}
其中显式 $\hbar$ 已由测度、分母和正则化尺度的幂次精确抵消。

用 Feynman 参数
\begin{equation*}
 \frac1{AB}=\int_0^1\frac{\dd x}{[xA+(1-x)B]^2}
\end{equation*}
并令 \(L=\ell+xp\)，分母化为
\begin{equation*}
 L^2-\Delta_s(x)+\ii0^+,
 \qquad
 \Delta_s(x)=m_R^2c^2-x(1-x)s.
\end{equation*}
于是
\begin{equation*}
 I(s)=\int_0^1\dd x\,
 \mu_{\rm DR}^{2\epsilon}
 \int\frac{\dd^d L}{(2\pi)^d}
 \frac{\ii}{(L^2-\Delta_s+\ii0^+)^2}.
\end{equation*}
做 Wick 转动 \(L^0=\ii L_E^0\)：
\begin{equation*}
 \dd^d L=\ii\dd^d L_E,
 \qquad
 L^2-\Delta+\ii0^+=-(L_E^2+\Delta-\ii0^+),
\end{equation*}
故
\begin{equation*}
 I(s)=-\int_0^1\dd x\,
 \mu_{\rm DR}^{2\epsilon}
 \int\frac{\dd^d L_E}{(2\pi)^d}
 \frac1{(L_E^2+\Delta_s-\ii0^+)^2}.
\end{equation*}
利用欧氏母积分
\begin{equation*}
 \int\frac{\dd^d L_E}{(2\pi)^d}
 \frac1{(L_E^2+\Delta)^2}
 =\frac1{(4\pi)^{d/2}}
 \Gamma\!\left(2-\frac d2\right)
 \Delta^{d/2-2},
\end{equation*}
代入 \(d=4-2\epsilon\)：
\begin{equation*}
 I(s)=-\int_0^1\dd x\,
 \frac1{(4\pi)^2}
 \Gamma(\epsilon)(4\pi)^\epsilon
 \left(\frac{\mu_{\rm DR}^2}{\Delta_s-\ii0^+}\right)^\epsilon.
\end{equation*}
逐项展开
\begin{align*}
 \Gamma(\epsilon)&=\frac1\epsilon-\gamma_E+\order(\epsilon),\\
 (4\pi)^\epsilon&=1+\epsilon\ln4\pi+\order(\epsilon^2),\\
 \left(\frac{\mu_{\rm DR}^2}{\Delta_s-\ii0^+}\right)^\epsilon
 &=1+\epsilon\ln\frac{\mu_{\rm DR}^2}{\Delta_s-\ii0^+}
 +\order(\epsilon^2).
\end{align*}
相乘到有限阶并执行参数积分，得到
\begin{equation*}
 I(s)=-\frac1{16\pi^2}
 \left[
 \frac1{\bar\epsilon}
 -\int_0^1\dd x\,
 \ln\frac{\Delta_s(x)-\ii0^+}{\mu_{\rm DR}^2}
 \right]+
 \order(\epsilon),
\end{equation*}
即\eqnrefs{eq:phi4-bubble-result-dp10}。当 \(\Delta_s(x)\) 穿过零时，\(-\ii0^+\) 决定对数的虚部，并与在壳中间态的 Cutkosky 不连续度一致。
\end{derivation}

\begin{derivation}[从\eqnrefs{eq:phi4-bare-coupling-r68}到\eqnrefs{eq:phi4-beta-one-loop-dp10}]
记

\begin{equation*}
 a_1\equiv\frac{3}{32\pi^2},
 \qquad
 P_\epsilon\equiv\frac1{\bar\epsilon}.
\end{equation*}
一圈阶裸耦合为
\begin{equation*}
 g_{4,B}
 =\mu^{2\epsilon}
 \left[g_4(\mu)+a_1P_\epsilon g_4^2(\mu)
 +\order(g_4^3)\right].
\end{equation*}
按正文\eqnrefs{eq:phi4-beta-one-loop-dp10} 中 $\beta_{g_4}$ 的定义，$d=4-2\epsilon$ 时固定裸耦合的总尺度导数写成
\begin{equation*}
 \mu\frac{\dd g_4}{\dd\mu}\Big|_{g_{4,B}}
 =-2\epsilon g_4+\beta_{g_4}(g_4).
\end{equation*}
对裸表达式求导时，反项同样通过 \(g_4(\mu)\) 依赖于 \(\mu\)：
\begin{align*}
 0={}&2\epsilon
 \left(g_4+a_1P_\epsilon g_4^2\right)\\
 &+\left(1+2a_1P_\epsilon g_4\right)
 \left(-2\epsilon g_4+\beta_{g_4}\right)
 +\order(g_4^3).
\end{align*}
保留到 \(g_4^2\) 后，\(2\epsilon g_4\) 项相消，并得到
\begin{equation*}
 0=\beta_{g_4}
 -2a_1\epsilon P_\epsilon g_4^2
 +\order(g_4^3,\epsilon g_4^2).
\end{equation*}
因为
\begin{equation*}
 \epsilon P_\epsilon
 =\epsilon\left(\frac1\epsilon-\gamma_E+\ln4\pi\right)
 \longrightarrow1,
\end{equation*}
四维极限给出
\begin{equation*}
 \beta_{g_4}=2a_1g_4^2+
\order(g_4^3)
 =\frac{3g_4^2}{16\pi^2}+
\order(g_4^3).
\end{equation*}
这里关键的因子 2 来自 \(d=4-2\epsilon\) 下的 \(\mu^{2\epsilon}\)；若把反项对 \(g_4(\mu)\) 的隐含尺度依赖漏掉，就会错误地得到相反符号。
\end{derivation}

\begin{derivation}[从\eqnrefs{eq:bare-ren-green-relation-r68}到\eqnrefs{eq:callan-symanzik-general-dp10}]

固定所有裸参数，对正文\eqnrefs{eq:bare-ren-green-relation-r68} 求 $\mu$ 的全导数：
\begin{equation*}
 0=\mu\frac{\dd}{\dd\mu}
 \left[Z_\varphi^{n/2}G_R^{(n)}(\mu,\{g_a\},\{m_b\})\right].
\end{equation*}
除以 $Z_\varphi^{n/2}$ 后，场强因子给出
\begin{equation*}
 \mu\frac{\dd}{\dd\mu}\ln Z_\varphi^{n/2}
 =n\gamma_\varphi.
\end{equation*}
而 Green 函数的全导数沿全部重整化参数方向展开为
\begin{equation*}
 \mu\frac{\dd G_R^{(n)}}{\dd\mu}
 =\left(
 \mu\partial_\mu
 +\sum_a\mu\frac{\dd g_a}{\dd\mu}\partial_{g_a}
 +\sum_b\mu\frac{\dd m_b}{\dd\mu}\partial_{m_b}
 \right)G_R^{(n)}.
\end{equation*}
代入正文\eqnrefs{eq:rg-general-flow-definitions-dp74} 的 $\beta$ 函数定义便得到正文\eqnrefs{eq:callan-symanzik-general-dp10}。单耦合、单质量理论只是把上述求和各保留一项后的特例。
\end{derivation}

\begin{derivation}[从\eqnrefs{eq:partial-wave-expansion-dp10}到\eqnrefs{eq:partial-wave-unitarity-bound-dp10}]
单通道弹性区间内

\begin{equation*}
 S_\ell=1+2\ii a_\ell,
 \qquad
 S_\ell S_\ell^*=1.
\end{equation*}
展开得到
\begin{align*}
 (1+2\ii a_\ell)(1-2\ii a_\ell^*)
 &=1+2\ii(a_\ell-a_\ell^*)+4|a_\ell|^2,\\
 &=1-4\Im a_\ell+4|a_\ell|^2.
\end{align*}
与 1 比较即有
\begin{equation*}
 \Im a_\ell=|a_\ell|^2.
\end{equation*}
写 \(a_\ell=x_\ell+\ii y_\ell\)，则
\begin{equation*}
 y_\ell=x_\ell^2+y_\ell^2,
\end{equation*}
配方得到
\begin{equation*}
 x_\ell^2+\left(y_\ell-\frac12\right)^2=\frac14.
\end{equation*}
因此精确弹性振幅位于以 \(\ii/2\) 为圆心、半径为 \(1/2\) 的 Argand 圆上。等价地，令
\begin{equation*}
 S_\ell=\ee^{2\ii\delta_\ell},
\end{equation*}
则
\begin{equation*}
 a_\ell
 =\frac{\ee^{2\ii\delta_\ell}-1}{2\ii}
 =\ee^{\ii\delta_\ell}\sin\delta_\ell,
\end{equation*}
从而
\begin{equation*}
 |a_\ell|\le1,
 \qquad
 |\Re a_\ell|\le\frac12.
\end{equation*}
树级振幅通常为实数，而精确弹性振幅必须位于上述 Argand 圆上。因此可采用保守的树级微扰一致性判据
\begin{equation*}
|\Re a_\ell^{\rm tree}|\le\frac12.
\end{equation*}
若树级结果越过该界，说明固定阶低能展开已经不能近似满足精确部分波幺正性，需要加入更高阶或新的有效自由度。
\end{derivation}



